\chapter{Реальный градуированный бимодуль и first-order квадрат}

Для \(A_F=\mathbb C\oplus\mathbb C\) неприводимые бимодули маркируются
парой индексов
\begin{equation}
 H_{ij},
 \qquad i,j\in\{X,C\},
\end{equation}
где левое действие использует \(i\)-ю компоненту алгебры, а правое ---
\(j\)-ю. Для матричного элемента
\(D_{ij,kl}:H_{ij}\to H_{kl}\) условие первого порядка имеет множитель
\begin{equation}
 (a_i-a_k)(b_j-b_l)D_{ij,kl}.
\end{equation}
Следовательно, ненулевое ребро разрешено только если \(i=k\) или \(j=l\):
оно может менять один индекс, но не оба одновременно.

\section{Минимальная reality-замкнутая диаграмма}

Прямая связь \(H_{XX}\leftrightarrow H_{CC}\) запрещена. Минимальный путь
\begin{equation}
 H_{XX}\longleftrightarrow H_{XC}
 \longleftrightarrow H_{CC}
\end{equation}
после замыкания относительно
\begin{equation}
 J:H_{ij}\longmapsto H_{ji}
\end{equation}
обязательно добавляет \(H_{CX}\). Получается квадрат
\begin{equation}
 H_{XX}\leftrightarrow H_{XC}\leftrightarrow H_{CC}
 \leftrightarrow H_{CX}\leftrightarrow H_{XX}.
\end{equation}

Градуировка
\begin{equation}
 \gamma=+1\quad\text{на }H_{XX}\oplus H_{CC},
 \qquad
 \gamma=-1\quad\text{на }H_{XC}\oplus H_{CX}
\end{equation}
делает все четыре разрешённых ребра нечётными.

\begin{theorem}[Минимальный first-order vertex pass]
Четырёхсекторный квадрат является минимальным \(J\)-замкнутым
\(A_F\)-бимодулем, который содержит оба диагональных факторных сектора и
допускает ненулевой нечётный оператор, удовлетворяющий first-order
правилу. Тем самым операторный запрет минимального диагонального
представления преодолён.
\end{theorem}

\section{Общий real-совместимый odd operator}

В порядке положительных секторов \((H_{XX},H_{CC})\) и отрицательных
\((H_{XC},H_{CX})\) нечётный самосопряжённый оператор имеет вид
\begin{equation}
 D_F=
 \begin{pmatrix}
 0&B\\
 B^\dagger&0
 \end{pmatrix},
 \qquad
 B=
 \begin{pmatrix}
 x&\overline x\\
 \overline z&z
 \end{pmatrix}.
\end{equation}
Здесь условие \(JD_F=D_FJ\) уже учтено. Оно оставляет два комплексных
параметра \(x,z\). Связующий блок невырожден тогда и только тогда, когда
\begin{equation}
 \det B=xz-\overline{xz}
 =2i\,\operatorname{Im}(xz)\ne0.
\end{equation}
Следовательно, вещественный оператор имеет ранг-дефицит; полный ранг требует
ненулевой относительной комплексной фазы.

\section{KO-размерностный red-team}

На четырёхсекторном квадрате транспонирование индексов сохраняет
диагональность и внедиагональность, поэтому
\begin{equation}
 J\gamma=\gamma J.
\end{equation}
Это не знак конечной KO-размерности \(6\), для которой требуется
\(J\gamma=-\gamma J\). При сохранении диагональных секторов минимальное
исправление состоит в particle--antiparticle doubling с противоположной
градуировкой, после которого \(J\) обменивает две копии. Размерность
конечного модуля возрастает с \(4\) до \(8\).

\begin{nogocriterion}[KO6 undoubled no-go]
Недиублированный четырёхсекторный квадрат допускает first-order vertex, но
не реализует знак \(J\gamma=-\gamma J\). Если Версия III требует конечную
KO-размерность \(6\), particle--antiparticle doubling является
обязательным, а правило физического half-trace должно быть выведено до
сравнения нормировок.
\end{nogocriterion}

\section{Параметрический статус}

Даже после KO-дублирования алгебра и first-order condition определяют
разрешённые позиции \(D_F\), но не значения \(x,z\). Минимальный квадрат
решает вопрос существования вершины, однако не фиксирует её модуль и
комплексную фазу. Поэтому он ещё не является parameter-free parent action.

\begin{protocol}[Следующий гейт]
Необходимо проверить, может ли автоморфизм обмена факторов, reality,
ориентируемость или спектральное действие:
\begin{enumerate}
 \item свести \(x,z\) к одной дискретно нормированной орбите;
 \item сохранить ненулевой \(\operatorname{Im}(xz)\);
 \item вывести physical half-trace после KO-дублирования;
 \item не изменить уже полученные факторные нормы.
\end{enumerate}
\end{protocol}

Машинный аудит сохранён в
\texttt{s2t\_v3\_real\_bimodule\_square\_results.json}.